\section{条件期望与线性回归是两种投影}
\label{sec:12-Real-Analysis-Support/07-Lp-Spaces-and-Function-Geometry/04}

一份回归报告摆在你面前。作者画了残差对每个特征的散点图，附上一行说明：残差与全部特征的相关系数都是 \(0\)（小数点后十五位），可见模型已经把特征里的信息榨干了。

你把 \(X\) 的平方项加进去重跑，\(R^2\) 从 \(0.31\) 跳到 \(0.79\)。

两件事并不矛盾，但作者读错了其中一件。残差与特征零相关不是模型好的证据——它是最小二乘这个算法必然产出的结果，无论模型对不对。要把这句话说清楚，得先承认线性回归和条件期望其实是同一个动作。

\subsection{把随机变量当成点，均方误差就是距离}

固定概率空间 \((\Omega,\mathcal F,P)\)，把平方可积的随机变量收进 \(L^2(P)\)，内积取

\[
\ip{U}{V}=\E[UV],\qquad \norm{U}_2=\sqrt{\E[U^2]} .
\]

这只是把上一节的 \(\int fg\,d\mu\) 换了个记号：随机变量本来就是 \(\Omega\) 上的可测函数，概率测度本来就是测度。换记号之后，一句统计学的话被翻译成了一句几何的话——\(\E[(Y-\widehat Y)^2]\) 是距离的平方，所以``最小化均方误差''就是``在某个集合里找离 \(Y\) 最近的点''。

集合是哪个，是唯一需要你拿主意的事。线性回归允许的预测是 \(X\) 的仿射函数，

\[
M_{\mathrm{lin}}=\set{a+b^{\trans}X:\ a\in\R,\ b\in\R^d} ,
\]

有限维，自动是闭子空间。条件期望允许的预测是 \(X\) 的任意平方可积可测函数，也就是 \(L^2(\sigma(X))\)，无限维，闭性来自 \(L^2\) 收敛保持可测性。两者嵌套：

\[
M_{\mathrm{lin}}\subseteq L^2(\sigma(X))\subseteq L^2(P) .
\]

以下把向子空间 \(M\) 的正交投影记作 \(\Pi_M\)（\(P\) 这个字母本节留给概率）。

\begin{center}
\begin{tikzpicture}[>=stealth]
% ---- 环境空间 ----
\draw[thick] (0.1,0.3) ellipse (4.6cm and 2.9cm);
\node at (-2.4,-1.75) {\footnotesize \(L^2(P)\)};
% ---- 大子空间：sigma(X) 可测函数 ----
\draw[thick] (-3.6,-0.9) -- (2.0,-1.9) -- (3.8,0.5) -- (-1.8,1.5) -- cycle;
\draw[thin] (2.75,0.70) -- (2.9,1.02);
\node at (2.9,1.18) {\footnotesize \(L^2(\sigma(X))\)};
% ---- 小子空间：仿射函数 ----
\draw[thick] (-2.2,-0.55) -- (1.2,-1.15) -- (1.75,-0.42) -- (-1.65,0.18) -- cycle;
\draw[thin] (1.75,-0.58) -- (1.58,-1.78);
\node at (1.30,-2.10) {\footnotesize \(\set{a+b^{\trans}X}\)};
% ---- 三个点 ----
\fill (-0.9,2.55) circle (0.06);
\node[above] at (-0.9,2.62) {\footnotesize \(Y\)};
\fill (-0.9,0.62) circle (0.06);
\node[anchor=west] at (-0.42,0.80) {\footnotesize \(\E[Y\given X]\)};
\fill (-0.4,-0.42) circle (0.06);
\node at (0.40,-0.68) {\footnotesize \(a+b^{\trans}X\)};
% ---- 两次投影 ----
\draw[thick] (-0.9,2.55) -- (-0.9,0.62);
\draw[thick] (-0.9,0.62) -- (-0.4,-0.42);
\draw[dashed] (-0.9,2.55) -- (-0.4,-0.42);
\node[left] at (-1.0,1.95) {\footnotesize \(Y-\E[Y\given X]\)};
\node[right] at (-0.64,1.66) {\footnotesize \(Y-a-b^{\trans}X\)};
\node[right] at (-0.25,0.24) {\footnotesize 直线够不着的那一截};
% ---- 两处直角标记，两小段线手绘 ----
\draw (-0.585,0.564) -- (-0.585,0.884);
\draw (-0.585,0.884) -- (-0.9,0.94);
\draw (-0.105,-0.472) -- (-0.235,-0.202);
\draw (-0.235,-0.202) -- (-0.53,-0.15);
\end{tikzpicture}
\end{center}

\emph{同一个 \(Y\) 向两个嵌套的子空间各投一次。虚线那条回归残差比实线那条长出来的部分，就是``\(X\) 能解释、直线解释不了''的那一截。}

\subsection{正规方程是正交条件的坐标写法}

先看小的那个平行四边形。投影定理说 \(Y-\Pi_{\mathrm{lin}}Y\) 与 \(M_{\mathrm{lin}}\) 中每个元素正交。子空间由 \(1,X_1,\dots,X_d\) 张成，逐个基取一遍就够：

\[
\E[Y-a-b^{\trans}X]=0,\qquad
\E\bigl[(Y-a-b^{\trans}X)X_j\bigr]=0,\quad j=1,\dots,d .
\]

\(d=1\) 时解出来是

\[
b=\frac{\Cov(X,Y)}{\Var(X)},\qquad a=\E Y-b\,\E X ,
\]

一元最小二乘的斜率与截距。\hyperref[sec:03-Linear-Algebra/03-Orthogonality/03]{``最小二乘：为无解方程寻找最佳近似''}里那条 \(A^{\trans}(y-A\hat x)=0\) 是同一句话的样本版本：把 \(\E\) 换成对样本求平均，内积换成 \(\R^n\) 里的点积，正交条件立刻变成正规方程。两处的几何一模一样，区别只在内积用的是总体分布还是经验分布。

再看大的那个。\(Y-\E[Y\given X]\) 与 \(L^2(\sigma(X))\) 中每个元素正交，也就是

\[
\E\bigl[(Y-\E[Y\given X])\,Z\bigr]=0\qquad\text{对所有 }Z\in L^2(\sigma(X)) .
\]

取 \(Z=\ind_G\)（\(G\in\sigma(X)\)），这一行退化成 \(\int_G Y\,dP=\int_G\E[Y\given X]\,dP\)，正是 \hyperref[sec:12-Real-Analysis-Support/05-Measure-Integration-and-Conditional-Expectation/05]{``条件期望是对信息 \(\sigma\)-代数的投影''}给出的定义性质。积分匹配是逐事件的说法，正交是几何的说法，在 \(L^2\) 中两者等价。塔式性质 \(\E[\E[Y\given\mathcal G]\given\mathcal H]=\E[Y\given\mathcal H]\)（\(\mathcal H\subseteq\mathcal G\)）也不再需要单独记忆：向大子空间投一次再向小子空间投一次，等于直接向小子空间投，这是嵌套投影的基本性质。

同一个定理，两个子空间。剩下的事情，是把这句话的后果一条条兑出来。

\subsection{线性回归拟合的是条件期望，不是数据}

第一条后果需要一句精确的表述。\(Y\) 拆成 \(\E[Y\given X]+\varepsilon\)，其中 \(\varepsilon\) 与整个 \(L^2(\sigma(X))\) 正交，当然也与更小的 \(M_{\mathrm{lin}}\) 正交，所以它向 \(M_{\mathrm{lin}}\) 的投影是零。投影算子是线性的，于是

\[
\Pi_{\mathrm{lin}}Y=\Pi_{\mathrm{lin}}\E[Y\given X] .
\]

翻译过来：总体最小二乘只看得见 \(\E[Y\given X]\) 这条曲线，看不见噪声。等价的写法更直观，

\[
(a,b)=\argmin_{a,b}\ \E\bigl[(Y-a-b^{\trans}X)^2\bigr]
=\argmin_{a,b}\ \E\bigl[(\E[Y\given X]-a-b^{\trans}X)^2\bigr] ,
\]

两式之差是一个与 \((a,b)\) 无关的常数 \(\E[\varepsilon^2]\)——Pythagoras 定理拆出来的那一项。所以``线性回归是条件期望的最佳线性近似''不是一句励志的类比，它字面成立：回归系数就是把 \(\E[Y\given X]\) 这条曲线在 \(X\) 的分布加权下做最小二乘拟合的结果。

这句话立刻带出一个容易踩的坑：加权用的是 \(X\) 的边缘分布。同一条 \(\E[Y\given X]\) 曲线，换一个采样方案、换一个总体、把 \(X\) 的取值范围截掉一段，拟合出的斜率就变。曲线是弯的时，回归系数不是可以搬来搬去的结构参数，它绑着你手上这批 \(X\) 的分布。\hyperref[sec:11-Mathematical-Statistics/08-Regression-and-ANOVA/01]{``经典线性回归把系数估计接到抽样分布''}讨论的是同一个系数在重复抽样下的波动，那是另一层不确定性；这里说的是即使样本无穷大、波动为零，系数仍随 \(X\) 的分布而动。

只有当 \(\E[Y\given X]\) 本来就是仿射的，两个平行四边形在这个方向上贴合，斜率才不再依赖 \(X\) 的分布。联合 Gaussian 是这种情形的标准例子：

\[
\E[Y\given X]=\E Y+\frac{\Cov(X,Y)}{\Var(X)}(X-\E X) .
\]

最佳预测与最佳线性预测重合。不是线性方法在这里变强了，是 Gaussian 世界里的条件结构本来就是线性的。离开这一族，图上那两个平行四边形重新分开。

\subsection{模型容量的几何含义就是子空间有多大}

第二条后果关于容量。设 \(M_1\subseteq M_2\) 是两个闭子空间，塔式性质给出 \(\Pi_1\Pi_2=\Pi_1\)，而 \(Y-\Pi_2Y\) 与 \(M_2\) 正交、\(\Pi_2Y-\Pi_1Y\) 落在 \(M_2\) 里，一次 Pythagoras 就得到

\[
\norm{Y-\Pi_1Y}_2^2=\norm{Y-\Pi_2Y}_2^2+\norm{\Pi_2Y-\Pi_1Y}_2^2 .
\]

子空间只要变大，误差就只减不增，多出来的那一块正好等于两次投影之差的长度平方——图上从 \(a+b^{\trans}X\) 指向 \(\E[Y\given X]\) 的那条短线段。

于是``加特征、加交互项、换成核方法或神经网络，训练误差单调下降''不必当成经验规律记：它是嵌套子空间上的一个恒等式，和数据长什么样无关。\(R^2\) 只增不减也是同一句话的另一种读法，这正是它作为模型比较指标不可靠的根源。\hyperref[sec:11-Mathematical-Statistics/08-Regression-and-ANOVA/02]{``ANOVA 用平方和分解比较多组均值''}里的平方和分解，本体就是上面这行等式在``全体均值''与``分组均值''两个嵌套子空间上的实例。

要留意这套几何全部发生在总体上：内积是 \(\E\)，不是样本平均。换成样本，内积变成经验分布下的平均，子空间变大照样让训练误差单调下降，但你想控制的是总体内积下的误差，两个内积之间的偏差不由投影几何管辖。\hyperref[sec:13-Machine-Learning/10-Learning-Theory/01]{``训练误差为何总是乐观''}处理的就是这条缝。几何告诉你容量大一定拟合得更好，它对容量大是否值得只字不提。

\subsection{残差与特征正交是定义性质，不是拟合好坏的证据}

第三条后果回到开头那份报告。正交条件 \(\E[(Y-a-b^{\trans}X)X_j]=0\) 是投影定理的结论，最小二乘解出来必然满足它——哪怕真实的 \(\E[Y\given X]\) 是一条抛物线、一个阶跃、一团乱麻。报告里那些十五位小数的零，度量的是求解器有没有把正规方程解对，不是模型有没有说对。

图上看得很清楚：那条虚线残差垂直于小平行四边形，是画法决定的；但它在大平行四边形里还剩着一整截，就是从 \(a+b^{\trans}X\) 指向 \(\E[Y\given X]\) 的那条线段。这一截是``\(X\) 本来能解释、而直线够不着''的部分，它按定义与 \(X\) 的一次项正交，所以在原始散点图上完全隐身。

可检验的东西因此只能是子空间之外的方向。把残差对 \(X_j^2\)、对交互项 \(X_jX_k\)、对拟合值本身、对时间顺序、对某个未进模型的分组变量作图——这些量都不在 \(M_{\mathrm{lin}}\) 里，正交条件管不着它们，残差在这些方向上一旦显出结构，就说明 \(\E[Y\given X]\) 确实有直线接不住的成分。开头那次 \(R^2\) 从 \(0.31\) 跳到 \(0.79\)，做的正是这件事：把子空间朝二次方向撑大一点，那一截原本藏起来的长度就被回收了。

\subsection{把角度换成别的，这套话就不成立了}

最后对条件过一遍账。投影身份要求 \(Y\in L^2\)：只有一阶矩时条件期望仍由积分匹配定义得好好的，但``最佳均方预测''这句话失去意义，因为均方本身不存在。子空间必须闭，\(M_{\mathrm{lin}}\) 靠有限维、\(L^2(\sigma(X))\) 靠可测性在 \(L^2\) 极限下保持——上一节那个稠密而不闭的多项式子空间提醒过，这一条不是走过场。

正交也别读成独立。线性回归的正交只给出 \(\E[\varepsilon]=0\) 与 \(\E[\varepsilon X_j]=0\)，是不相关；条件期望的正交给出 \(\E[\varepsilon\given X]=0\)，是均值独立，比不相关强，但仍弱于独立——残差的方差、偏度完全可以随 \(X\) 变化。三个层级依次收紧，混用它们是异方差和条件分布误判的常见来源。

还有一层边界在损失函数上。整套推理建立在均方即距离平方之上；换成绝对误差，你落回 \(L^1\)，上一节算过那里没有内积，条件中位数与分位数回归都不是投影，它们的``残差正交''版本根本写不出来。这也解释了为什么分位数回归的估计要靠线性规划而不是解一组线性方程。

一个单元到此走完了它的四步：先给函数装上距离，让``两条曲线差多少''成为一个有答案的问题；再要求完备，让逼近序列的极限不逃出空间；接着在 \(p=2\) 处装上角度，让``最近点''有了唯一解和可检验的判据；最后发现统计里那些看起来各不相干的最优——最佳线性预测、最佳均方预测、方差分解、Fourier 截断——是同一个投影动作落在不同子空间上的结果。再往前，泛函分析要问对偶空间、算子谱与弱拓扑，那已越出本书为概率与学习准备的地基范围。
